\documentclass[12pt]{article}

\usepackage[margin=1in]{geometry}
\usepackage{setspace}
\usepackage{graphicx}
\usepackage{hyperref}
\usepackage{enumitem}

\title{Course Pathway and Instructor-Quality Visualizer for UofT CS}
\author{James Wong \and Tommy Le \and Jerry Ruoheng Ma}
\date{\today}

\begin{document}
\maketitle
\onehalfspacing

\section{Introduction}
\textbf{Project Question/Goal:} How can we help UofT Computer Science students plan prerequisite pathways to target courses while incorporating instructor quality information from RateMyProfessors?

Our project models CS courses and prerequisite relations as a directed graph and augments this graph with instructor rating signals. The motivation is that existing planning tools show prerequisite rules, but do not combine them with teaching-quality data in one workflow. We aim to support planning decisions by combining (1) prerequisite reachability and (2) course-level instructor rating summaries.

\section{Datasets}
\subsection{Dataset 1: UofT CS Course Catalog}
\begin{itemize}[leftmargin=*]
    \item Name: \texttt{Datasets/CourseData.csv}
    \item Source: UofT Arts \& Science Academic Calendar (scraped in the reference project and reused here).
    \item Format: pipe-delimited CSV.
    \item Fields used by program: course code, title, description, prerequisites, recommended, corequisite, exclusion, breadth requirement.
    \item Notes: The file has been moved to the project root and is loaded directly by \texttt{course\_dataset.py}.
\end{itemize}

\subsection{Dataset 2: RateMyProf Course--Professor Ratings}
\begin{itemize}[leftmargin=*]
    \item Name: \texttt{Datasets/CourseProfessorRatings.csv}
    \item Source: RateMyProfessors, scraped from UofT St.\ George teacher data.
    \item Format: CSV with columns:
    \begin{itemize}[leftmargin=*]
        \item \texttt{course\_number}
        \item \texttt{professors\_by\_score\_json}
    \end{itemize}
    \item Data shape used by program: \texttt{course\_number -> \{score: [professor\_names]\}}.
    \item Notes: Scores are 1--5 (stored as decimals such as 4.2). This dataset is generated by our scraper pipeline and then loaded as a local dataset.
\end{itemize}

\section{Computational Overview}
\subsection{Data Representation (Trees/Graphs)}
Current implementation includes typed data models for courses and course-rating summaries:
\begin{itemize}[leftmargin=*]
    \item \texttt{Course} in \texttt{models.py}
    \item \texttt{ProfessorProfile} in \texttt{models.py}
    \item \texttt{CourseProfessorRatings} in \texttt{models.py}
\end{itemize}
The current computation phase already builds a directed prerequisite graph from \texttt{Datasets/CourseData.csv}. Nodes are normalized course numbers and directed edges encode prerequisite constraints.

\subsection{Major Computations and Algorithms}
Implemented so far:
\begin{itemize}[leftmargin=*]
    \item Course catalog loading and normalization (\texttt{course\_dataset.py})
    \item RateMyProf scraping (\texttt{ratemyprof\_scraper.py}) using paginated school-level teacher search and course-code extraction
    \item Course-rating dataset generation and loading (\texttt{rmp\_course\_dataset.py})
    \item Directed prerequisite graph construction and basic recommendation logic (\texttt{prerequisite\_graph.py})
    \item Website route and interactive graph rendering (\texttt{web\_app.py})
\end{itemize}
The scraper collects UofT St.\ George professors in Computer Science, extracts each professor's course codes and average rating, and aggregates to per-course score buckets.

\subsection{Output, Visualization, and Interactivity}
Current output includes both terminal summaries and a working website baseline. In the website users:
\begin{itemize}[leftmargin=*]
    \item input completed courses and an optional target course;
    \item view a Plotly prerequisite graph with color-coded course states;
    \item view simple ranked next-course recommendations using prerequisite unlock count and average instructor rating.
\end{itemize}

\subsection{External Python Libraries}
Current implementation uses:
\begin{itemize}[leftmargin=*]
    \item \texttt{networkx} for graph algorithms and graph construction,
    \item \texttt{plotly} for interactive graph visualization,
    \item \texttt{flask} for the website backend.
\end{itemize}

\section{Instructions: Dataset Access and Running the Program}
\subsection{Python and Library Setup}
\begin{enumerate}[leftmargin=*]
    \item Install Python 3.13.5.
    \item Install dependencies from \texttt{requirements.txt} (to be finalized as project modules stabilize).
\end{enumerate}

\subsection{Dataset Download Instructions}
\begin{itemize}[leftmargin=*]
    \item \texttt{Datasets/CourseData.csv} is included in the project.
    \item \texttt{Datasets/CourseProfessorRatings.csv} can be regenerated from live scraping or included in the submission zip.
\end{itemize}

\subsection{How to Run}
\begin{enumerate}[leftmargin=*]
    \item Run \texttt{python main.py} to verify local dataset loading.
    \item To build ratings dataset from live RateMyProf data, run:
    \texttt{python -c "from main import run\_build\_rmp\_dataset; run\_build\_rmp\_dataset()"}
    \item To load and inspect generated ratings dataset, run:
    \texttt{python -c "from main import run\_load\_rmp\_dataset\_summary; run\_load\_rmp\_dataset\_summary()"}
    \item To run the website, run:
    \texttt{python -c "from main import run\_web\_app; run\_web\_app()"}
\end{enumerate}

\section{Changes from Proposal to Final Submission}
Compared to the proposal, our implemented data approach is now concrete:
\begin{itemize}[leftmargin=*]
    \item We moved course data into a local dataset folder (\texttt{Datasets/CourseData.csv}).
    \item We implemented a reproducible RateMyProf scraping pipeline that outputs a second local CSV dataset.
    \item We standardized rating aggregation as a course-number dictionary keyed by score, matching our recommender input design.
    \item We implemented an initial website and prerequisite graph view rather than only terminal output.
\end{itemize}

\section{Discussion}
So far, the main result is a working two-dataset pipeline plus an interactive website baseline. The website can render the prerequisite graph, mark completed/unlocked/target-path nodes, and compute first-pass course recommendations. A key limitation is that scraped coverage depends on publicly available RateMyProf entries and may miss some course-professor combinations. We handle this by keeping all course numbers in the dataset, with empty mappings for courses without discovered rating data.

Next steps are:
\begin{itemize}[leftmargin=*]
    \item improve prerequisite parsing for complex ``and/or`` prerequisite expressions,
    \item upgrade recommendation scoring with stronger weighting and constraints,
    \item expand website controls (filters, views, and pathway comparison).
\end{itemize}

\section{References}
\begin{thebibliography}{9}
\bibitem{rmp}
RateMyProfessors. \url{https://www.ratemyprofessors.com/}

\bibitem{uoftcalendar}
University of Toronto Arts \& Science Academic Calendar (Computer Science section). \url{https://artsci.calendar.utoronto.ca/section/Computer-Science}

\bibitem{networkx}
NetworkX Documentation. \url{https://networkx.org/}

\bibitem{plotly}
Plotly Python Documentation. \url{https://plotly.com/python/}

\bibitem{flask}
Flask Documentation. \url{https://flask.palletsprojects.com/}
\end{thebibliography}

\end{document}
